\documentclass[titlepage,fontsize=10.5pt]{jlreq}

%---If you use beamer class, add aspectratio=169 and change jlreq to beamer---%

%---Package part---%
\usepackage{amssymb,amsmath,amsthm,amscd}
\usepackage{ascmac}
\usepackage{bm}
\usepackage{color}
\usepackage{derivative}
\usepackage{enumitem} % If you use beamer metropolis, ignore this package
\usepackage{fancybox}
\usepackage{frame,framed}
\usepackage{float}
\usepackage{graphicx}
\usepackage{mathtools}
\usepackage{physics2}
\usepackage{siunitx}
\usepackage{tabularray}
\usepackage{tcolorbox}
\usepackage{tikz}
\usepackage{url}
\usepackage{titlepic} % If you use beamer metropolis, ignore this package
\usepackage{wrapfig}
%---Package part END---%

%---Package Module part---%
\usephysicsmodule{ab,braket,doubleprod,xmat}
%---Package Module part END---%

%---Cont foot nad head settings---%
\DefTblrTemplate{contfoot-text}{normal}{（次ページに続く）}
\SetTblrTemplate{contfoot-text}{normal}
\DefTblrTemplate{conthead-text}{normal}{（続き）}
\SetTblrTemplate{conthead-text}{normal}
%---Cont foot nad head settings END---%

%---For equition number part---%
\usepackage[hang,small,bf]{caption}
\usepackage[subrefformat=parens]{subcaption}
\captionsetup{compatibility=false}
\renewcommand{\theequation}{\thesection.\arabic{equation}}
\makeatletter
\@addtoreset{equation}{section}
\makeatother
%---For equition number part END---%

%---For beamer settings part---%
%---If you use beamer metropolis,remove header % from below sentense---%
% \usepackage[haranoaji,deluxe]{luatexja-preset}% Font setting
% \usepackage{luatexja}% For use japanese
% \renewcommand{\kanjifamilydefault}{\gtdefault}% Setting normal font to gothic
% \usetheme[block=fill]{metropolis}% For use metropolis theme
%---For beamer settings part END---%

\begin{document}

\title{電子顕微鏡実習}
\author{平山大知}
\date{\today}
\maketitle

%---For make index part---%
\tableofcontents
\clearpage
%---For make index part END---%

\section{目的} % 目的(fold)
    \label{sec.目的}
    走査電子顕微鏡について理解を深めるために、実際の試料の製作および観察を通して、その動作原理と基本的な操作方法を習得すること。
    %目的(end)
\section{背景原理} %背景原理(fold)
    \label{sec.背景原理}
    
    \subsection{電子線} % (fold)
    \label{sub:電子線}
    電子顕微鏡は試料を観察するために電子線を用いる。電子線の波長$\lambda$は以下の式で表される。
    \begin{align}
        \lambda = \dfrac{h}{\sqrt{2meV}}
    \end{align}
    ここに$h$はプランク定数、$m$は電子の質量、$e$は電子の電荷、$V$は電子の加速電圧である。ここに電子の質量と電荷とプランク定数を代入すると以下の式を得る。
    \begin{align}
        \lambda \,[\unit{\nano\metre}]= \dfrac{1.23}{\sqrt{V\,[\unit{\volt}]}}
    \end{align}
    例えば加速電圧が\qty{20}{\kilo\volt}であるならば、電子線の波長は\qty{8.69e-3}{\nano\metre}となる。つまり電子を\qty{20}{\kilo\volt}で加速する電子顕微鏡は$\qty{8.69e-3}{\nano\metre}$の分解能をもつ。光学式顕微鏡であれば光の波長、およそ\qty{400}{\nano\metre}の分解能しか持たないわけであるから、電子顕微鏡は光学式顕微鏡に対して約$46000$倍の分解能を持つことになる。

    式からわかるように単に加速電圧を上げていけば分解能も上がるが、先の式は古典論の範囲であるから、相対論的効果が現れるほど加速した場合には修正が必要である。
    % subsection 電子線 (end)

    \subsection{走査型電子顕微鏡} % (fold)
    \label{sub:走査型電子顕微鏡}
    電子顕微鏡には走査型電子顕微鏡と透過型電子顕微鏡がある。今回使用したのは走査型電子顕微鏡であるためこれについて記述する。

    走査型電子顕微鏡では試料上で電子線照射を走査して、そのときに出てくる二次電子や反射電子、特性X線等の信号を観測することで試料を調べる。

    二次電子は加速して照射した電子が表面の原子の価電子を放出させたものであり、表面の情報を多く含んでいる。なので二次電子から得られる像は表面の微細な構造の情報をもつ。

    反射電子は加速して照射した電子が原子にあたって反射されたものであり、内部についての情報を多く含んでいる。反射電子の起源が原子の反射によるから、原子についての情報を持つことになる。

    また特性X線の情報を得ることで構成元素の解析も可能となる。

    電子線が原子内部を通過する際に、原子内の電子と衝突し、電子が弾かれた場合弾かれた電子を補うように外殻から電子が移動してくる。そのときポテンシャルが下がっているのでその差分をX線として放出する。このポテンシャルの差分は原子の種類と移動した殻の種類により決まるため、このX線を見ることで原子の特性を観測できる。

    原子核と電子線の相互作用で失われるエネルギーもあり、これもまたX線として放出される。このX線は連続X線と呼ばれてバックグラウンドの信号として得られる。

    \subsubsection{試料作成} % (fold)
    \label{ssub:試料作成}
    
    % subsubsection 試料作成 (end)
    
    % subsection 走査型電子顕微鏡 (end)

    %背景原理(end)
\section{実験手段・機材}
    \label{sec.実験手段・機材}
    \subsection{試料の作成} % (fold)
    \label{sub:試料の作成}
    髪の毛を実際に観察するために髪の毛を用いて試料作成を行った。

    試料台にカーボン製の両面テープを貼付し、その上に髪の毛(\qty{2}{\centi\metre}ほどを)貼付した。

    その後にオートファインコータに試料台を設置し、オートファインコータの真空を引いた後にスパッタリング電流\qty{30}{\milli\volt}に設定、\qty{25}{\sec}コーティングを行う設定として試料台上の試料へコーディングを行った。
    % subsection 試料の作成 (end)

    \subsection{試料観察} % (fold)
    \label{sub:試料観察}
    髪の毛、キャピラリープレート、10円玉を観察した。
    % subsection 試料観察 (end)

    \subsection{構成元素の解析} % (fold)
    \label{sub:構成元素の解析}
    キャピラリープレート、ICチップ部品、鉱石について構成元素の解析を行った。
    % subsection 構成元素の解析 (end)
\section{結果}

\section{考察}

\section{まとめ}



\end{document}